\section{Conclusion}
\label{sec:conclusion}

\subsection{Summary}

Cellula meets every requirement of the subject: a recursive array whose dimension is chosen at
run time, neighborhoods exposed as traversals and definable by the user, the eleven operators
organized in a class hierarchy, an input file format, a graphical display for one, two and more
dimensions, a menu of the available files, the two required definitions, and an answer to the
bonus question obtained with the program itself.

Three design decisions proved the most valuable.

\begin{itemize}
    \item \emph{Letting the grid answer for every position.} Because \code{Grid.cellAt} always
          returns a cell, neither the neighborhoods nor the operators contain a single bounds check,
          and a second boundary condition was added by writing one method.
    \item \emph{Generating the neighborhoods instead of listing them.} Two formulas replace six
          special cases and make the program work in four dimensions without a line of
          dimension-specific code.
    \item \emph{Separating the definition from the automaton, and the automaton from its display.}
          The same classes serve the window, the terminal, the batch commands and the tests; this is
          what made it possible to test nearly everything automatically and to produce the figures of
          this report with the program.
\end{itemize}

\subsection{Limits}

\begin{itemize}
    \item \textbf{Repeated work in rules.} A rule such as the Game of Life mentions
          \code{COMPTER(G8*)} three times, and each occurrence walks the neighborhood again. Sharing
          identical sub-expressions when the tree is built would roughly halve the cost of such rules.
    \item \textbf{Memory.} Every cell is an object holding its coordinates. This is what the
          recursive, object-oriented design of the subject leads to, and it limits practical grids to
          a few million cells; a flat array of integers would use ten times less memory.
    \item \textbf{Display of higher dimensions.} A grid of three dimensions or more is only seen
          one plane at a time. A true 3D view would be a natural addition.
    \item \textbf{Cycle detection} is bounded to periods of at most 1\,024 generations, and relies
          on 64-bit fingerprints: a collision is extremely unlikely but not impossible.
    \item \textbf{The classification of pictures} in the rule explorer is a heuristic. It
          shortlists candidates well, but a human still chooses among them.
\end{itemize}

\subsection{Possible extensions}

The architecture leaves clear places for further work: new operators are one class and one line
of registration (\code{MOD}, or a count of cells in an arbitrary state, would make multi-state
automata such as Wireworld concise); a new boundary condition is one constant; another input
format only has to produce an \code{AutomatonDefinition}; and a new display is another
implementation of \code{Projection} or another listener of the simulation.
